% =============================================================================
% 063 / 068
% Status: raw
% =============================================================================
% Notes:
%
% Source question:
% 推主可以用吳語播音體翻譯一遍東亞語言對比材料中的名篇《北風與太陽》嗎？這篇小故事覆蓋了好多語言現象。
% 
% 附原文：“ 有一回，北風跟太陽在那裡爭，誰的本事大。爭來爭去分不出高低。這時候路上來了個過路人，身上穿著件厚大衣。他們倆就講好了，誰有本事先叫這個過路人把他的大衣脫下來，就算誰的本事大。北風就拼命地吹起來，不過他越是吹得厲害，那個過路人就把大衣包得越緊。後來北風沒辦法了，只好拉倒了。過了一會兒，太陽出來了。他熱烘烘地一曬，那個過路人馬上就把那件大衣脫下來了。這下北風只好承認，他們倆當中還是太陽的本事大。”
% =============================================================================

\chapter[推主可以用吳語播音體翻譯一遍東亞語言對比材料中的名篇《北風與太陽》嗎？這篇小故事覆蓋了好多語言現象。 附原文：“ 有一回，北風跟太陽在那裡爭，誰的本事大。爭來爭去分不出高低。這時候路上來了個過路人，身上穿著件厚大衣。他們倆就講好了，誰有本事先叫這個過路人把他的大衣脫下來，就算誰的本事大。北風就拼命地吹起來，不過他越是吹得厲害，那個過路人就把大衣包得越緊。後來北風沒辦法了，只好拉倒了。過了一會兒，太陽出來了。他熱烘烘地一曬，那個過路人馬上就把那件大衣脫下來了。這下北風只好承認，他們倆當中還是太陽的本事大。”]{推主可以用吳語播音體翻譯一遍東亞語言對比材料中的名篇《北風與太陽》嗎？這篇小故事覆蓋了好多語言現象。 \\[0.35em] 附原文：“ 有一回，北風跟太陽在那裡爭，誰的本事大。爭來爭去分不出高低。這時候路上來了個過路人，身上穿著件厚大衣。他們倆就講好了，誰有本事先叫這個過路人把他的大衣脫下來，就算誰的本事大。北風就拼命地吹起來，不過他越是吹得厲害，那個過路人就把大衣包得越緊。後來北風沒辦法了，只好拉倒了。過了一會兒，太陽出來了。他熱烘烘地一曬，那個過路人馬上就把那件大衣脫下來了。這下北風只好承認，他們倆當中還是太陽的本事大。”}
\qaanswer{
可以試一試
}

